\documentclass[a4paper,12pt]{article}

\input{../../Preambulo.tex}
\newcommand{\assunto}{CONJUNTOS E FUNÇÕES}
\newcommand{\atividade}{lista} % prova ou lista
\newcommand{\turma}{ELETROELETRÔNICA 1.\textordmasculine\ ANO}
\newcommand{\numProva}{A}
\newcommand{\professor}{Igor Martins Silva}
\newcommand{\bimestre}{3}
\newcommand{\ano}{2026}
\newcommand{\mesTexto}{agosto}
\newcommand{\mesNum}{08}
\newcommand{\dia}{27}
\newcommand{\dataTexto}{\dia\ de \mesTexto\ de \ano}
\newcommand{\dataNun}{\dia/\mesNum/\ano}

\newcommand{\titulo}{}

\ifdefstring{\atividade}{prova}
  {\renewcommand{\titulo}{Avaliação de Matemática}}
{
\ifdefstring{\atividade}{lista}
  {\renewcommand{\titulo}{Lista de exercícios de Matemática}}
{
\ifdefstring{\atividade}{as}
  {\renewcommand{\titulo}{Avaliação Somativa}}
{
\ifdefstring{\atividade}{ad}
  {\renewcommand{\titulo}{Avaliação Diagnóstica}}
  {\renewcommand{\titulo}{Atividade de Matemática}}
}}}

\newcommand{\emtipo}{%
  \ifdefstring{\tipo}{individual}{\tipo}{em \tipo}%
}

\edef\pathCapa{\currfiledir}

\newcommand{\subtitulobase}[3]{%
    % #1 = título da 1ª coluna
    % #2 = conteúdo da 1ª coluna
    % #3 = mostra número da prova? (vazio ou algo)
    
    \noindent
    \begin{minipage}{2.8cm}
        \includegraphics[width=2.8cm]{\pathCapa Logo_Cefet.png}
    \end{minipage}
    \hfill
    \begin{minipage}{\dimexpr\textwidth-6.2cm\relax}
        \begin{center}
        CEFET-MG -- Campus Contagem \\[3pt]
        \titulo\ -- \bimestre.\textordmasculine\ bimestre de \ano \\[3pt]
        \professor
    \end{center}
    \end{minipage}
    \hfill
    \begin{minipage}{3.2cm}
        \includegraphics[width=3.2cm]{\pathCapa Brasao_Brasil.png}
    \end{minipage}
    
    \begin{center}
        \vspace{15pt}
        \begin{tabular}{|C{210pt}|C{70pt}|C{210pt}|}
            \hline
            \textbf{#1} &
            \textbf{DATA} &
            \textbf{TURMA} \\
            \hline
            #2 & \dataNun & \turma \\
            \hline
        \end{tabular}
    \end{center}
    
    #3
    
    \vspace{15pt}
}

\newcommand{\subtitulolis}{%
  \subtitulobase{ASSUNTO}{\assunto}{}%
}

\newcommand{\subtitulopr}{%
  \subtitulobase{ASSUNTO}{\assunto}{\boxed{\numProva}}%
}

\newcommand{\subtituloas}{%
  \subtitulobase{DISCIPLINA}{MATEMÁTICA}{\boxed{\numProva}}%
}

\newcommand{\subtituload}{%
  \subtitulobase{DISCIPLINA}{MATEMÁTICA}{}%
}

\newcommand{\valorquestao}{}

\newcommand{\definevalorquestao}{%
    \ifdefstring{\atividade}{lista}{%
        % Não faz nada para lista
    }{%
        \ifdefstring{\modo}{objetivo}{%
            \renewcommand{\valorquestao}{\ValorQObj ponto}
        }{%
            \renewcommand{\valorquestao}{\ValorQDisc pontos}
        }%
    }%
}

\newenvironment{exercicioBanco}[1][]{%
    \ifdefstring{\atividade}{lista}{%
        \begin{exercicio}%
    }{%
        \begin{exercicio}[#1]%
    }%
}{%
    \end{exercicio}
}

\ifdefstring{\atividade}{lista}{
    \pagecolor{background-color}
}{
    
}



\hypersetup{
    pdfauthor={\professor},
    pdftitle={\titulo \ -- \assunto \ -- \turma},
}

\title{\titulo \ -- \assunto \ -- \turma}
\author{\professor}
\date{\data}

\graphicspath{
    {../../Conjuntos/}
}

%************* INÍCIO DO DOCUMENTO *************
\begin{document}
    \subtitulolis
        
    \phantomsection
    \addcontentsline{toc}{section}{Exercício 1}
    %%%%%%%%%%%%%%%%%%%%%%%%%%%%%%%%%
%
%       EXERCÍCIO
%
%%%%%%%%%%%%%%%%%%%%%%%%%%%%%%%%%

\ifdefstring{\atividade}{prova}{%
    \ifdefstring{\modo}{objetivo}{%
        \renewcommand{\valorquestao}{\ValorQObj\ ponto}
    }{%
        \renewcommand{\valorquestao}{\ValorQDisc\ pontos}
    }%
}%

\begin{exercicioBanco}[\valorquestao]
Julgue as afirmações a seguir, classificando cada uma como verdadeira (V) ou falsa (F).
%
\begin{enumerate}[\(\bullet\)]
    \item \(\{0,3,5\} \subseteq \{0,1,2,3,4\}\).
    %
    \item O conjunto \(\{x \mid x\) é vogal da palavra ``estacionamento''\(\}\) é igual ao conjunto \(\{a, e, i, o\}\).
    %
    \item \(\varnothing \in \bbN\).
    %
    \item \(\{\text{Lua}\} \subseteq \{x \mid x\) é satélite natural da Terra\(\}\).
\end{enumerate}
%

% Define as alternativas
\newcommand{\alternativas}{%
    \begin{center}
        \begin{tabularx}{\textwidth}{XXXXX}
            (a) (F, V, F, V). &
            (b) (V, F, V, F). &
            (c) (F, V, V, F). &
            (d) (V, F, F, V). &
            (e) (F, F, V, F).
        \end{tabularx}
    \end{center}
}

% Define a resposta correta
\newcommand{\resposta}{A}

% Lógica condicional para exibição
\ifdefstring{\atividade}{lista}{%
    \alternativas
    \vspace{0.5em}
    
    \noindent\textbf{Resposta:} letra \textbf{\resposta}.
}{%
    \ifdefstring{\modo}{objetivo}{%
        \alternativas
    }{%
        % modo = discursiva → não mostra alternativas
    }
}
\end{exercicioBanco}


    \noindent\rule{\linewidth}{1pt}  % linha horizontal fina
    
    \phantomsection
    \addcontentsline{toc}{section}{Exercício 2}
    %%%%%%%%%%%%%%%%%%%%%%%%%%%%%%%%%
%
%       EXERCÍCIO
%
%%%%%%%%%%%%%%%%%%%%%%%%%%%%%%%%%

\ifdefstring{\atividade}{prova}{%
    \ifdefstring{\modo}{objetivo}{%
        \renewcommand{\valorquestao}{\ValorQObj\ ponto}
    }{%
        \renewcommand{\valorquestao}{\ValorQDisc\ pontos}
    }%
}%

\begin{exercicioBanco}[\valorquestao]
Dado o conjunto \(A = \{\{1\}, 2, 3\}\), assinale a alternativa correta.

% Define as alternativas
\newcommand{\alternativas}{%
    \begin{center}
        \begin{tabularx}{\textwidth}{XXXXX}
            (a) \(1 \in A\). &
            (b) \(1 \subseteq A\). &
            (c) \(\{1\} \subseteq A\). &
            (d) \(\{1\} \in A\). &
            (e) \(\varnothing \in A\).
        \end{tabularx}
    \end{center}
}

% Define a resposta correta
\newcommand{\resposta}{D}

% Lógica condicional para exibição
\ifdefstring{\atividade}{lista}{%
    \alternativas
    \vspace{0.5em}
    
    \noindent\textbf{Resposta:} letra \textbf{\resposta}.
}{%
    \ifdefstring{\modo}{objetivo}{%
        \alternativas
    }{%
        % modo = discursiva → não mostra alternativas
    }
}
\end{exercicioBanco}


    \noindent\rule{\linewidth}{1pt}  % linha horizontal fina
    
    \phantomsection
    \addcontentsline{toc}{section}{Exercício 3}
    %%%%%%%%%%%%%%%%%%%%%%%%%%%%%%%%%
%
%       EXERCÍCIO
%
%%%%%%%%%%%%%%%%%%%%%%%%%%%%%%%%%

\ifdefstring{\atividade}{prova}{%
    \ifdefstring{\modo}{objetivo}{%
        \renewcommand{\valorquestao}{\ValorQObj\ ponto}
    }{%
        \renewcommand{\valorquestao}{\ValorQDisc\ pontos}
    }%
}%

\begin{exercicioBanco}[\valorquestao]
No diagrama a seguir, \(A, B\) e \(C\) são três conjuntos não vazios.

\begin{center}
    \begin{tikzpicture}[scale=0.5]
        % Conjunto B (maior)
        \draw[thick, purple!70] (0,0) circle (3);

        % Conjunto A
        \draw[thick, green!60!black] (-1.2,0) circle (1);

        % Conjunto C
        \draw[thick, orange!90!black] (1.5,0) circle (1);

        % Rótulos com bolinhas coloridas
        \node[circle, fill=purple!70, text=black, inner sep=1pt] at (-1.8,2.5) {\small$B$};
        \node[circle, fill=green!60!black, text=black, inner sep=1pt] at (-2,0.8) {\small$A$};
        \node[circle, fill=orange!90!black, text=black, inner sep=1pt] at (2.2,0.8) {\small$C$};
    \end{tikzpicture}
\end{center}

Assinale verdadeira (V) ou falso (F) a cada uma das seguintes sentenças.
%
\begin{enumerate}[I.]
    \begin{multicols}{5}
    \item \(A \cap B = \varnothing\).
    %
    \item \(A \subseteq B\).
    %
    \item \(\varnothing \subseteq A\).
    %
    \item \(A \cup C = B\).
    %
    \item \(B\setminus A = C\).
    \end{multicols}
\end{enumerate}

% Define as alternativas
\newcommand{\alternativas}{%
    \begin{center}
        \begin{tabularx}{\textwidth}{XXXXX}
            (a) (F, V, F, V, V). &
            (b) (V, V, V, F, F). &
            (c) (F, V, F, V, F). &
            (d) (V, F, V, F, F). &
            (e) (V, F, V, F, V).
        \end{tabularx}
    \end{center}
}

% Define a resposta correta
\newcommand{\resposta}{B}

% Lógica condicional para exibição
\ifdefstring{\atividade}{lista}{%
    \alternativas
    \vspace{0.5em}
    
    \noindent\textbf{Resposta:} letra \textbf{\resposta}.
}{%
    \ifdefstring{\modo}{objetivo}{%
        \alternativas
    }{%
        % modo = discursiva → não mostra alternativas
    }
}
\end{exercicioBanco}


    \noindent\rule{\linewidth}{1pt}  % linha horizontal fina
    
    \phantomsection
    \addcontentsline{toc}{section}{Exercício 4}
    %%%%%%%%%%%%%%%%%%%%%%%%%%%%%%%%%
%
%       EXERCÍCIO
%
%%%%%%%%%%%%%%%%%%%%%%%%%%%%%%%%%

\ifdefstring{\atividade}{prova}{%
    \ifdefstring{\modo}{objetivo}{%
        \renewcommand{\valorquestao}{\ValorQObj\ ponto}
    }{%
        \renewcommand{\valorquestao}{\ValorQDisc\ pontos}
    }%
}%

\begin{exercicioBanco}[\valorquestao]
Em uma festa, $35$ pessoas discutiam sobre dois filmes, $A$ e $B$. 
Sabe-se que:
%
\begin{enumerate}[\(\bullet\)]
    \begin{multicols}{2}
        \item $15$ pessoas assistiram ao filme $A$;
        \item $8$ pessoas assistiram aos dois filmes;
        \item $10$ pessoas não assistiram a nenhum dos dois filmes.
    \end{multicols}
\end{enumerate}
%
Sabendo que todas as $35$ pessoas opinaram, o número de pessoas que assistiram ao filme $B$ é:


% Define as alternativas
\newcommand{\alternativas}{%
    \begin{center}
        \begin{tabularx}{\textwidth}{XXXXX}
            (a) \(7\). &
            (b) \(8\). &
            (c) \(10\). &
            (d) \(18\). &
            (e) \(25\).
        \end{tabularx}
    \end{center}
}

% Define a resposta correta
\newcommand{\resposta}{D}

% Lógica condicional para exibição
\ifdefstring{\atividade}{lista}{%
    \alternativas
    \vspace{0.5em}
    
    \noindent\textbf{Resposta:} letra \textbf{\resposta}.
}{%
    \ifdefstring{\modo}{objetivo}{%
        \alternativas
    }{%
        % modo = discursiva → não mostra alternativas
    }
}
\end{exercicioBanco}


    \noindent\rule{\linewidth}{1pt}  % linha horizontal fina
    
    \phantomsection
    \addcontentsline{toc}{section}{Exercício 5}
    %%%%%%%%%%%%%%%%%%%%%%%%%%%%%%%%%
%
%       EXERCÍCIO
%
%%%%%%%%%%%%%%%%%%%%%%%%%%%%%%%%%

\ifdefstring{\atividade}{prova}{%
    \ifdefstring{\modo}{objetivo}{%
        \renewcommand{\valorquestao}{\ValorQObj\ ponto}
    }{%
        \renewcommand{\valorquestao}{\ValorQDisc\ pontos}
    }%
}%

\begin{exercicioBanco}[\valorquestao]
Seja \(M(a)\) o conjunto dos múltiplos de \(a\) e \(D(a)\) o conjunto dos divisores de \(a\), com \(a \in \bbN = \{1, 2, 3, \ldots\}\). Liste os elementos do conjunto \(M(3) \cap D(30)\).

% Define as alternativas
\newcommand{\alternativas}{%
    \begin{center}
        \begin{tabularx}{\textwidth}{XXX}
            (a) \(\{6,15,30\}\). &
            (b) \(\{1,2,3,5,6,10,15,30\}\). &
            (c) \(\{3,6,15,30\}\). \\[5pt]
            (d) \(\{3,15\}\). &
            (e) \(\{3,6,9,12,15,18,21,24,27,30\}\).
        \end{tabularx}
    \end{center}
}

% Define a resposta correta
\newcommand{\resposta}{C}

% Lógica condicional para exibição
\ifdefstring{\atividade}{lista}{%
    \alternativas
    \vspace{0.5em}
    
    \noindent\textbf{Resposta:} letra \textbf{\resposta}.
}{%
    \ifdefstring{\modo}{objetivo}{%
        \alternativas
    }{%
        % modo = discursiva → não mostra alternativas
    }
}
\end{exercicioBanco}


    \noindent\rule{\linewidth}{1pt}  % linha horizontal fina
    
    \phantomsection
    \addcontentsline{toc}{section}{Exercício 6}
    %%%%%%%%%%%%%%%%%%%%%%%%%%%%%%%%%
%
%       EXERCÍCIO
%
%%%%%%%%%%%%%%%%%%%%%%%%%%%%%%%%%

\ifdefstring{\atividade}{prova}{%
    \ifdefstring{\modo}{objetivo}{%
        \renewcommand{\valorquestao}{\ValorQObj\ ponto}
    }{%
        \renewcommand{\valorquestao}{\ValorQDisc\ pontos}
    }%
}%

\begin{exercicioBanco}[\valorquestao]
No universo $\bbN = \{1,2,3,\ldots\}$, sejam:
%
\begin{enumerate}[\(\bullet\)]
    \begin{multicols}{2}
        \item $A$ o conjunto dos números pares;
        \item $B$ o conjunto dos números múltiplos de $3$;
        \item $C$ o conjunto dos números múltiplos de $5$.
    \end{multicols}
\end{enumerate}
%
Determine os $3$ menores números que pertencem ao conjunto
%
\[
    B\setminus(A \cup C).
\]
%

% Define as alternativas
\newcommand{\alternativas}{%
    \begin{center}
        \begin{tabularx}{\textwidth}{XXXXX}
            (a) \(3, 6, 9\). &
            (b) \(3, 9, 21\). &
            (c) \(3, 6, 15\). &
            (d) \(6, 9, 12\). &
            (e) \(6, 9, 15\).
        \end{tabularx}
    \end{center}
}

% Define a resposta correta
\newcommand{\resposta}{B}

% Lógica condicional para exibição
\ifdefstring{\atividade}{lista}{%
    \alternativas
    \vspace{0.5em}
    
    \noindent\textbf{Resposta:} letra \textbf{\resposta}.
}{%
    \ifdefstring{\modo}{objetivo}{%
        \alternativas
    }{%
        % modo = discursiva → não mostra alternativas
    }
}
\end{exercicioBanco}


    \noindent\rule{\linewidth}{1pt}  % linha horizontal fina
    
    \phantomsection
    \addcontentsline{toc}{section}{Exercício 7}
    %%%%%%%%%%%%%%%%%%%%%%%%%%%%%%%%%
%
%       EXERCÍCIO
%
%%%%%%%%%%%%%%%%%%%%%%%%%%%%%%%%%

\ifdefstring{\atividade}{prova}{%
    \ifdefstring{\modo}{objetivo}{%
        \renewcommand{\valorquestao}{\ValorQObj\ ponto}
    }{%
        \renewcommand{\valorquestao}{\ValorQDisc\ pontos}
    }%
}%

\begin{exercicioBanco}[\valorquestao]
Os números \(x\) e \(y\) são tais que \(5 \le x \le 10\) e \(20 \le y \le 30\). Determine o maior valor possível de \(\dfrac{x}{y}\).

% Define as alternativas
\newcommand{\alternativas}{%
    \begin{center}
        \begin{tabularx}{\textwidth}{XXXXX}
            (a) \(\dfrac{1}{6}\). &
            (b) \(\dfrac{1}{4}\). &
            (c) \(\dfrac{1}{3}\). &
            (d) \(\dfrac{1}{2}\). &
            (e) \(1\).
        \end{tabularx}
    \end{center}
}

% Define a resposta correta
\newcommand{\resposta}{D}

% Lógica condicional para exibição
\ifdefstring{\atividade}{lista}{%
    \alternativas
    \vspace{0.5em}
    
    \noindent\textbf{Resposta:} letra \textbf{\resposta}.
}{%
    \ifdefstring{\modo}{objetivo}{%
        \alternativas
    }{%
        % modo = discursiva → não mostra alternativas
    }
}
\end{exercicioBanco}


    \noindent\rule{\linewidth}{1pt}  % linha horizontal fina
    
    \phantomsection
    \addcontentsline{toc}{section}{Exercício 8}
    %%%%%%%%%%%%%%%%%%%%%%%%%%%%%%%%%
%
%       EXERCÍCIO
%
%%%%%%%%%%%%%%%%%%%%%%%%%%%%%%%%%

\ifdefstring{\atividade}{prova}{%
    \ifdefstring{\modo}{objetivo}{%
        \renewcommand{\valorquestao}{\ValorQObj\ ponto}
    }{%
        \renewcommand{\valorquestao}{\ValorQDisc\ pontos}
    }%
}%

\begin{exercicioBanco}[\valorquestao]
Dados os conjuntos \(A = [-1,6[\), \(B = ]-4,2]\) e \(E = ]-2,4[\), calcule \(E\setminus(A\cap B)\).

% Define as alternativas
\newcommand{\alternativas}{%
    \begin{center}
        \begin{tabularx}{\textwidth}{XXX}
            (a) \(]-2,-1[ \cup ]2,4[\). &
            (b) \(]-2,-1] \cup [2,4[\). &
            (c) \(]-2,2]\). \\[5pt]
            (d) \([-1,2]\). &
            (e) \(]-2,4[ \setminus [-1,2[\).
        \end{tabularx}
    \end{center}
}

% Define a resposta correta
\newcommand{\resposta}{A}

% Lógica condicional para exibição
\ifdefstring{\atividade}{lista}{%
    \alternativas
    \vspace{0.5em}
    
    \noindent\textbf{Resposta:} letra \textbf{\resposta}.
}{%
    \ifdefstring{\modo}{objetivo}{%
        \alternativas
    }{%
        % modo = discursiva → não mostra alternativas
    }
}
\end{exercicioBanco}


    \noindent\rule{\linewidth}{1pt}  % linha horizontal fina
    
    \phantomsection
    \addcontentsline{toc}{section}{Exercício 9}
    %%%%%%%%%%%%%%%%%%%%%%%%%%%%%%%%%
%
%       EXERCÍCIO
%
%%%%%%%%%%%%%%%%%%%%%%%%%%%%%%%%%

\ifdefstring{\atividade}{prova}{%
    \ifdefstring{\modo}{objetivo}{%
        \renewcommand{\valorquestao}{\ValorQObj\ ponto}
    }{%
        \renewcommand{\valorquestao}{\ValorQDisc\ pontos}
    }%
}%

\begin{exercicioBanco}[\valorquestao]
Assinale a afirmação verdadeira.

% Define as alternativas
\newcommand{\alternativas}{%
    \begin{center}
        \begin{tabularx}{\textwidth}{X}
            (a) \((\sqrt{3}+1)(\sqrt{3}-1)\) é irracional e \(0,999\cdots\) é racional. \\[5pt]
            (b) \((\sqrt{3}+1)(\sqrt{3}-1)\) é racional e \(0,999\cdots\) é racional. \\[5pt]
            (c) \((\sqrt{3}+1)(\sqrt{3}-1)\) é racional e \(0,999\cdots\) é irracional. \\[5pt]
            (d) \((\sqrt{3}+1)(\sqrt{3}-1)\) é irracional e \(0,999\cdots\) é irracional. \\[5pt]
            (e) \((\sqrt{3}+1)(\sqrt{3}-1)\) e \(0,999\cdots\) não são reais.
        \end{tabularx}
    \end{center}
}

% Define a resposta correta
\newcommand{\resposta}{B}

% Lógica condicional para exibição
\ifdefstring{\atividade}{lista}{%
    \alternativas
    \vspace{0.5em}
    
    \noindent\textbf{Resposta:} letra \textbf{\resposta}.
}{%
    \ifdefstring{\modo}{objetivo}{%
        \alternativas
    }{%
        % modo = discursiva → não mostra alternativas
    }
}
\end{exercicioBanco}


    \noindent\rule{\linewidth}{1pt}  % linha horizontal fina
    
    \phantomsection
    \addcontentsline{toc}{section}{Exercício 10}
    %%%%%%%%%%%%%%%%%%%%%%%%%%%%%%%%%
%
%       EXERCÍCIO
%
%%%%%%%%%%%%%%%%%%%%%%%%%%%%%%%%%

\ifdefstring{\atividade}{prova}{%
    \ifdefstring{\modo}{objetivo}{%
        \renewcommand{\valorquestao}{\ValorQObj\ ponto}
    }{%
        \renewcommand{\valorquestao}{\ValorQDisc\ pontos}
    }%
}%

\begin{exercicioBanco}[\valorquestao]
Sejam \(A = \{x \in \bbR \mid -4 \le x \le 3\}\) e \(B = \{x \in \bbR \mid -2 \le x < 5\}\). Determine \(A\setminus B\).

% Define as alternativas
\newcommand{\alternativas}{%
    \begin{center}
        \begin{tabularx}{\textwidth}{XXX}
            (a) \(\{x \in \bbR \mid 3 < x < 5\}\). &
            (b) \(\{x \in \bbR \mid -4 \le x \le -2\}\). &
            (c) \(\{x \in \bbR \mid -4 \le x < -2\}\). \\[5pt]
            (d) \(\{x \in \bbR \mid 3 \le x \le 5\}\). &
            (e) \(\{x \in \bbR \mid -2 \le x < 5\}\).
        \end{tabularx}
    \end{center}
}

% Define a resposta correta
\newcommand{\resposta}{C}

% Lógica condicional para exibição
\ifdefstring{\atividade}{lista}{%
    \alternativas
    \vspace{0.5em}
    
    \noindent\textbf{Resposta:} letra \textbf{\resposta}.
}{%
    \ifdefstring{\modo}{objetivo}{%
        \alternativas
    }{%
        % modo = discursiva → não mostra alternativas
    }
}
\end{exercicioBanco}


    \noindent\rule{\linewidth}{1pt}  % linha horizontal fina
    
    \phantomsection
    \addcontentsline{toc}{section}{Exercício 11}
    %%%%%%%%%%%%%%%%%%%%%%%%%%%%%%%%%
%
%       EXERCÍCIO
%
%%%%%%%%%%%%%%%%%%%%%%%%%%%%%%%%%

\ifdefstring{\atividade}{prova}{%
    \ifdefstring{\modo}{objetivo}{%
        \renewcommand{\valorquestao}{\ValorQObj\ ponto}
    }{%
        \renewcommand{\valorquestao}{\ValorQDisc\ pontos}
    }%
}%

\begin{exercicioBanco}[\valorquestao]
Considere a reta numérica.
%
\vspace{15pt}
\begin{center}
    \begin{tikzpicture}[scale=1]

    % Linha principal com seta
    \draw[->, thick] (-5.5,0) -- (5.5,0);

    % Marcas e números
    \foreach \x in {-5,-4,-3,-2,-1,0,1,2,3,4,5}
    {
        \draw (\x,0.2) -- (\x,-0.2);
        \node[below] at (\x,-0.2) {\small $\x$};
    }
    \end{tikzpicture}
\end{center}
%
O número \(-\dfrac{17}{5}\) foi marcado entre que pontos dessa reta numérica?

% Define as alternativas
\newcommand{\alternativas}{%
    \begin{center}
        \begin{tabularx}{\textwidth}{XXXXX}
            (a) \(2\) e \(3\). &
            (b) \(3\) e \(4\). &
            (c) \(-3\) e \(-2\). &
            (d) \(-4\) e \(-3\). &
            (e) \(-5\) e \(-4\).
        \end{tabularx}
    \end{center}
}

% Define a resposta correta
\newcommand{\resposta}{D}

% Lógica condicional para exibição
\ifdefstring{\atividade}{lista}{%
    \alternativas
    \vspace{0.5em}
    
    \noindent\textbf{Resposta:} letra \textbf{\resposta}.
}{%
    \ifdefstring{\modo}{objetivo}{%
        \alternativas
    }{%
        % modo = discursiva → não mostra alternativas
    }
}
\end{exercicioBanco}


    \noindent\rule{\linewidth}{1pt}  % linha horizontal fina
    
    \phantomsection
    \addcontentsline{toc}{section}{Exercício 12}
    %%%%%%%%%%%%%%%%%%%%%%%%%%%%%%%%%
%
%       EXERCÍCIO
%
%%%%%%%%%%%%%%%%%%%%%%%%%%%%%%%%%

\ifdefstring{\atividade}{prova}{%
    \ifdefstring{\modo}{objetivo}{%
        \renewcommand{\valorquestao}{\ValorQObj\ ponto}
    }{%
        \renewcommand{\valorquestao}{\ValorQDisc\ pontos}
    }%
}%

\begin{exercicioBanco}[\valorquestao]
Ordene os números racionais \(p = \dfrac{13}{24}\), \(q = \dfrac{2}{3}\) e \(r = \dfrac{5}{8}\).

% Define as alternativas
\newcommand{\alternativas}{%
    \begin{center}
        \begin{tabularx}{\textwidth}{XXXXX}
            (a) \(p < q < r\). &
            (b) \(r < p < q\). &
            (c) \(q < r < p\). &
            (d) \(r < q < p\). &
            (e) \(p < r < q\).
        \end{tabularx}
    \end{center}
}

% Define a resposta correta
\newcommand{\resposta}{E}

% Lógica condicional para exibição
\ifdefstring{\atividade}{lista}{%
    \alternativas
    \vspace{0.5em}
    
    \noindent\textbf{Resposta:} letra \textbf{\resposta}.
}{%
    \ifdefstring{\modo}{objetivo}{%
        \alternativas
    }{%
        % modo = discursiva → não mostra alternativas
    }
}
\end{exercicioBanco}


    \noindent\rule{\linewidth}{1pt}  % linha horizontal fina
    
    \phantomsection
    \addcontentsline{toc}{section}{Exercício 13}
    \input{../C_Exercicio13_1}
    \noindent\rule{\linewidth}{1pt}  % linha horizontal fina
    
    \phantomsection
    \addcontentsline{toc}{section}{Exercício 14}
    %%%%%%%%%%%%%%%%%%%%%%%%%%%%%%%%%
%
%       EXERCÍCIO
%
%%%%%%%%%%%%%%%%%%%%%%%%%%%%%%%%%

\ifdefstring{\atividade}{prova}{%
    \ifdefstring{\modo}{objetivo}{%
        \renewcommand{\valorquestao}{\ValorQObj\ ponto}
    }{%
        \renewcommand{\valorquestao}{\ValorQDisc\ pontos}
    }%
}%

\begin{exercicioBanco}[\valorquestao]
Seja \(f: \bbR \to \bbR\) uma função tal que \(f(\frac{1}{2}) = \sqrt{a}\), com \(a\) um número real positivo, e \(f(x+1) = xf(x)\), para todo \(x \in \bbR\). Determine \(f(\frac{7}{2})\).

% Define as alternativas
\newcommand{\alternativas}{%
    \begin{center}
        \begin{tabularx}{\textwidth}{XXXXX}
            (a) \(\dfrac{\sqrt{a}}{4}\). &
            (b) \(\dfrac{3\sqrt{a}}{8}\). &
            (c) \(\dfrac{9\sqrt{a}}{4}\). &
            (d) \(\dfrac{15\sqrt{a}}{8}\). &
            (e) \(\dfrac{21\sqrt{a}}{8}\).
        \end{tabularx}
    \end{center}
}

% Define a resposta correta
\newcommand{\resposta}{D}

% Lógica condicional para exibição
\ifdefstring{\atividade}{lista}{%
    \alternativas
    \vspace{0.5em}
    
    \noindent\textbf{Resposta:} letra \textbf{\resposta}.
}{%
    \ifdefstring{\modo}{objetivo}{%
        \alternativas
    }{%
        % modo = discursiva → não mostra alternativas
    }
}
\end{exercicioBanco}


    \noindent\rule{\linewidth}{1pt}  % linha horizontal fina
    
    \phantomsection
    \addcontentsline{toc}{section}{Exercício 15}
    %%%%%%%%%%%%%%%%%%%%%%%%%%%%%%%%%
%
%       EXERCÍCIO
%
%%%%%%%%%%%%%%%%%%%%%%%%%%%%%%%%%

\ifdefstring{\atividade}{prova}{%
    \ifdefstring{\modo}{objetivo}{%
        \renewcommand{\valorquestao}{\ValorQObj\ ponto}
    }{%
        \renewcommand{\valorquestao}{\ValorQDisc\ pontos}
    }%
}%

\begin{exercicioBanco}[\valorquestao]
Sejam \(A = \{0,1,2,3\}\) e \(B = \{x \in \bbZ \mid -6 \le x \le 6\}\). Considere a função
%
\[
    \begin{array}{rcl}
        f: \; A & \to & B \\
        x & \mapsto & 2x.
    \end{array}
\]
%
Determine o conjunto imagem de \(f\).

% Define as alternativas
\newcommand{\alternativas}{%
    \begin{center}
        \begin{tabularx}{\textwidth}{XXX}
            (a) \(B\). &
            (b) \(\{0,2,4,6\}\). &
            (c) \(A\). \\[5pt]
            (d) \(\bbZ\). &
            (e) \(\{x \in \bbZ \mid 0 \le x \le 6\}\).
        \end{tabularx}
    \end{center}
}

% Define a resposta correta
\newcommand{\resposta}{B}

% Lógica condicional para exibição
\ifdefstring{\atividade}{lista}{%
    \alternativas
    \vspace{0.5em}
    
    \noindent\textbf{Resposta:} letra \textbf{\resposta}.
}{%
    \ifdefstring{\modo}{objetivo}{%
        \alternativas
    }{%
        % modo = discursiva → não mostra alternativas
    }
}
\end{exercicioBanco}


    \noindent\rule{\linewidth}{1pt}  % linha horizontal fina
    
    \phantomsection
    \addcontentsline{toc}{section}{Exercício 16}
    %%%%%%%%%%%%%%%%%%%%%%%%%%%%%%%%%
%
%       EXERCÍCIO
%
%%%%%%%%%%%%%%%%%%%%%%%%%%%%%%%%%

\ifdefstring{\atividade}{prova}{%
    \ifdefstring{\modo}{objetivo}{%
        \renewcommand{\valorquestao}{\ValorQObj\ ponto}
    }{%
        \renewcommand{\valorquestao}{\ValorQDisc\ pontos}
    }%
}%

\begin{exercicioBanco}[\valorquestao]
Considere a função \(f:[-3,6] \to [-3,3]\), cujo gráfico está abaixo.
%
\begin{center}
    \begin{tikzpicture}[scale=0.4]

    % Eixos
    \draw[->] (-4,0) -- (7,0);      % eixo x
    \draw[->] (0,-4) -- (0,5);      % eixo y

    % Marcas principais
    \foreach \x in {-3,-2,-1,1,2,3,4,5,6}
        \draw (\x,0.15) -- (\x,-0.15);

    \foreach \y in {-3,-2,-1,1,2,3,4}
        \draw (0.15,\y) -- (-0.15,\y);

    % Rótulos importantes
    \node[below] at (-3,0) {$-3$};
    \node[below] at (2,0) {$2$};
    \node[right] at (0,1) {$1$};
    \node[below] at (6,0) {$6$};
    \node[left] at (0,3) {$3$};
    \node[left] at (0,-3) {$-3$};

    % Segmento inclinado
    \draw[thick] (-3,1) -- (2,-3);

    % Pontos fechados do segmento inclinado
    \fill (-3,1) circle (4pt);
    \fill (2,-3) circle (4pt);

    % Segmento horizontal superior
    \draw[thick] (2,3) -- (6,3);

    % Círculo aberto em (2,3)
    \draw[thick] (2,3) circle (4pt);

    % Ponto fechado em (6,3)
    \fill (6,3) circle (4pt);

    % Linhas tracejadas auxiliares
    \draw[dashed] (2,0) -- (2,3);
    \draw[dashed] (6,0) -- (6,3);
    \draw[dashed] (0,3) -- (2,3);
    \draw[dashed] (0,-3) -- (2,-3) -- (2,-1);
    \draw[dashed] (-3,0) -- (-3,1) -- (0,1);

    \end{tikzpicture}
\end{center}
%
Determine o valor de \(\dfrac{f(5)}{f(-3)-f(2)}\).

% Define as alternativas
\newcommand{\alternativas}{%
    \begin{center}
        \begin{tabularx}{\textwidth}{XXXXX}
            (a) \(\dfrac{1}{4}\). &
            (b) \(\dfrac{2}{4}\). &
            (c) \(\dfrac{3}{4}\). &
            (d) \(1\). &
            (e) \(\dfrac{5}{4}\).
        \end{tabularx}
    \end{center}
}

% Define a resposta correta
\newcommand{\resposta}{C}

% Lógica condicional para exibição
\ifdefstring{\atividade}{lista}{%
    \alternativas
    \vspace{0.5em}
    
    \noindent\textbf{Resposta:} letra \textbf{\resposta}.
}{%
    \ifdefstring{\modo}{objetivo}{%
        \alternativas
    }{%
        % modo = discursiva → não mostra alternativas
    }
}
\end{exercicioBanco}


    \noindent\rule{\linewidth}{1pt}  % linha horizontal fina
    
    \phantomsection
    \addcontentsline{toc}{section}{Exercício 17}
    %%%%%%%%%%%%%%%%%%%%%%%%%%%%%%%%%
%
%       EXERCÍCIO
%
%%%%%%%%%%%%%%%%%%%%%%%%%%%%%%%%%

\ifdefstring{\atividade}{prova}{%
    \ifdefstring{\modo}{objetivo}{%
        \renewcommand{\valorquestao}{\ValorQObj\ ponto}
    }{%
        \renewcommand{\valorquestao}{\ValorQDisc\ pontos}
    }%
}%

\begin{exercicioBanco}[\valorquestao]
Sejam os conjuntos \(A = \{1,2\}\) e \(B = \{0,1,2\}\). Assinale a alternativa cuja relação define uma função de \(A\) para \(B\).

% Define as alternativas
\newcommand{\alternativas}{%
    \begin{center}
        \begin{tabularx}{\textwidth}{XXX}
            (a) \(f(x) = 2x\). &
            (b) \(f(x) = x+1\). &
            (c) \(f(x) = x^{2} - 3x + 3\). \\[5pt]
            (d) \(f(x) = x^{2} - 2x\). &
            (e) \(f(x) = x - 2\).
        \end{tabularx}
    \end{center}
}

% Define a resposta correta
\newcommand{\resposta}{C}

% Lógica condicional para exibição
\ifdefstring{\atividade}{lista}{%
    \alternativas
    \vspace{0.5em}
    
    \noindent\textbf{Resposta:} letra \textbf{\resposta}.
}{%
    \ifdefstring{\modo}{objetivo}{%
        \alternativas
    }{%
        % modo = discursiva → não mostra alternativas
    }
}
\end{exercicioBanco}


    \noindent\rule{\linewidth}{1pt}  % linha horizontal fina
    
    \phantomsection
    \addcontentsline{toc}{section}{Exercício 18}
    %%%%%%%%%%%%%%%%%%%%%%%%%%%%%%%%%
%
%       EXERCÍCIO
%
%%%%%%%%%%%%%%%%%%%%%%%%%%%%%%%%%

\ifdefstring{\atividade}{prova}{%
    \ifdefstring{\modo}{objetivo}{%
        \renewcommand{\valorquestao}{\ValorQObj\ ponto}
    }{%
        \renewcommand{\valorquestao}{\ValorQDisc\ pontos}
    }%
}%

\begin{exercicioBanco}[\valorquestao]
Das figuras a seguir, indique aquela que representa o gráfico de uma função real \(f: [a,b] \to \bbR\), onde \(a, b \in \bbR\).

% Define as alternativas
\newcommand{\alternativas}{%
    \begin{center}
        \begin{tabularx}{\textwidth}{XXX}
            \hspace{40pt} (a) & \hspace{40pt} (b) & \hspace{40pt} (c) \\[5pt]
            \begin{tikzpicture}[scale=0.6]
                \draw[->,thick] (0,0)--(4,0) node[right]{$x$};
                \draw[->,thick] (0,0)--(0,4) node[above]{$y$};
                \draw[dashed] (1,0)--(1,1.1);
                \draw[dashed] (3.5,0)--(3.5,2.5);
                \draw (1,1.2)--(2.5,3)--(1.8,1.5)--(3.5,2.5);
                \draw node[below, yshift=-3pt] at (1,0) {\(a\)};
                \draw node[below] at (3.5,0) {\(b\)};
            \end{tikzpicture}
            &
            \begin{tikzpicture}[scale=0.6]
                \draw[->,thick] (0,0)--(4,0) node[right]{$x$};
                \draw[->,thick] (0,0)--(0,4) node[above]{$y$};
                \draw[dashed] (1,0)--(1,1);
                \draw[dashed] (3.5,0)--(3.5,2);
                \draw (1,1)--(2,1.5);
                \draw[dashed] (2.6,0)--(2.6,3);
                \draw (2.6,3)--(3.5,2);
                \draw[dashed] (2,0)--(2,1.5);
                \draw node[below, yshift=-3pt] at (1,0) {\(a\)};
                \draw node[below] at (3.5,0) {\(b\)};
            \end{tikzpicture}
            &
            \begin{tikzpicture}[scale=0.6]
                \draw[->,thick] (0,0)--(4,0) node[right]{$x$};
                \draw[->,thick] (0,0)--(0,4) node[above]{$y$};
                \draw[dashed] (1,0)--(1,3);
                \draw[dashed] (3.5,0)--(3.5,2.5);
                \draw (1,3)--(3.5,1);
                \draw (1,1.5)--(3.5,2.5);
                \draw node[below, yshift=-3pt] at (1,0) {\(a\)};
                \draw node[below] at (3.5,0) {\(b\)};
            \end{tikzpicture}
            \\[13pt]
            \hspace{40pt} (d) & \hspace{40pt} (e) & \\[5pt]
            \begin{tikzpicture}[scale=0.6]
                \draw[->,thick] (0,0)--(4,0) node[right]{$x$};
                \draw[->,thick] (0,0)--(0,4) node[above]{$y$};
                \draw[dashed] (1,0)--(1,1);
                \draw[dashed] (3.5,0)--(3.5,3);
                \draw (1,1)--(1.5,0.8)--(2,2)--(2.5,1)--(3.5,3);
                \draw node[below, yshift=-3pt] at (1,0) {\(a\)};
                \draw node[below] at (3.5,0) {\(b\)};
            \end{tikzpicture}
            &
            \begin{tikzpicture}[scale=0.6]
                \draw[->,thick] (0,0)--(4,0) node[right]{$x$};
                \draw[->,thick] (0,0)--(0,4) node[above]{$y$};
                \draw[dashed] (1,0)--(1,2);
                \draw[dashed] (3.5,0)--(3.5,2);
                \draw (2.25,2) ellipse (1.25 and 0.6);
                \draw node[below, yshift=-3pt] at (1,0) {\(a\)};
                \draw node[below] at (3.5,0) {\(b\)};
            \end{tikzpicture}
            &
        \end{tabularx}
    \end{center}
}

% Define a resposta correta
\newcommand{\resposta}{D}

% Lógica condicional para exibição
\ifdefstring{\atividade}{lista}{%
    \alternativas
    \vspace{0.5em}
    
    \noindent\textbf{Resposta:} letra \textbf{\resposta}.
}{%
    \ifdefstring{\modo}{objetivo}{%
        \alternativas
    }{%
        % modo = discursiva → não mostra alternativas
    }
}
\end{exercicioBanco}


\end{document}
%*************** FINAL DO DOCUMENTO ***************
